% LaTeX skabelon, master fil
% ------------------------------------------------------------------------------
% Dette er hovedfilen i projektet, hvori indholdet fra alle input-filer samles;
% tekst, billeder, literaturhenvisninger, osv.


% LaTeX' 'book' class er den documentclass, der har de mest fleksible muligheder
% for at projekt som dette. Se https://tex.stackexchange.com/a/36989/118167
%
% De ekstra argumenter i firkantklammerne sætter grundskriftstørrelsen til 11pt,
% sidestørrelse til A4, layout til to-sidet, og gør så nye chapters starter på
% en højreside.
\documentclass[11pt,a4paper,twoside,openright]{book}

% Variable til automatisk at sætte titler, forfattere, osv. diverse steder
% (såsom forsiden og titelbladet). For flere muligheder skal man justere de
% steder, de bliver brugt, manuelt.
\def \projecttitle       {\LaTeX{} Template}
\def \projectsubtitle    {A Helping Hand For Getting Started}
\def \projectdegree      {Mathematics-Economics}
\def \projectsemester    {1}
\def \projectsupervisors {Janus S. Valberg-Madsen}
\def \projectgroup       {Group XYZ}
\def \projectauthors     {
  Nova B. Valberg-Madsen\\
  Line B. Olsen
}

% Den medfølgende 'macros' fil definerer kommandoer, der genererer en titelside
% i AAU-stil.  Bemærk, at kommandoer normalt puttes ind i preamblet, men
% 'macros' er skrevet af en anden person, så derfor har jeg beholdt den som
% seperat fil.
%  A simple AAU report template.
%  2013-03-06 v. 1.0.0
%  Copyright 2010-2013 by Jesper Kjær Nielsen <jkn@es.aau.dk>
%
%  This is free software: you can redistribute it and/or modify
%  it under the terms of the GNU General Public License as published by
%  the Free Software Foundation, either version 3 of the License, or
%  (at your option) any later version.
%
%  This is distributed in the hope that it will be useful,
%  but WITHOUT ANY WARRANTY; without even the implied warranty of
%  MERCHANTABILITY or FITNESS FOR A PARTICULAR PURPOSE.  See the
%  GNU General Public License for more details.
%
%  You can find the GNU General Public License at <http://www.gnu.org/licenses/>.
%
%
%
% see, e.g., http://en.wikibooks.org/wiki/LaTeX/Customizing_LaTeX#New_commands
% for more information on how to create macros

%%%%%%%%%%%%%%%%%%%%%%%%%%%%%%%%%%%%%%%%%%%%%%%%
% Macros for the titlepage
%%%%%%%%%%%%%%%%%%%%%%%%%%%%%%%%%%%%%%%%%%%%%%%%
%Creates the aau titlepage
\newcommand{\aautitlepage}[3]{%
  {
    %set up various length
    \ifx\titlepageleftcolumnwidth\undefined
      \newlength{\titlepageleftcolumnwidth}
      \newlength{\titlepagerightcolumnwidth}
    \fi
    \setlength{\titlepageleftcolumnwidth}{0.5\textwidth-\tabcolsep}
    \setlength{\titlepagerightcolumnwidth}{\textwidth-2\tabcolsep-\titlepageleftcolumnwidth}
    %create title page
    \thispagestyle{empty}
    \noindent%
    \begin{tabular}{@{}ll@{}}
      \parbox{\titlepageleftcolumnwidth}{
        \includegraphics[width=\titlepageleftcolumnwidth]{./fig/img/aau_logo_en.pdf}
      } &
      \parbox{\titlepagerightcolumnwidth}{\raggedleft\sf\small
        #2
      }\bigskip\\
       #1 &
      \parbox[t]{\titlepagerightcolumnwidth}{%
      \textbf{Abstract:}\bigskip\par
        \fbox{\parbox{\titlepagerightcolumnwidth-2\fboxsep-2\fboxrule}{%
          #3
        }}
      }\\
    \end{tabular}
    \vfill
    \noindent{\footnotesize\emph{The content of this report is freely available, but publication (with reference) may only be pursued due to agreement with the authors.}}
    \clearpage
  }
}

%Create english project info
\newcommand{\englishprojectinfo}[8]{%
  \parbox[t]{\titlepageleftcolumnwidth}{
    \textbf{Title:}\\ #1\bigskip\par
    \textbf{Theme:}\\ #2\bigskip\par
    \textbf{Project Period:}\\ #3\bigskip\par
    \textbf{Project Group:}\\ #4\bigskip\par
    \textbf{Participant(s):}\\ #5\bigskip\par
    \textbf{Supervisor(s):}\\ #6\bigskip\par
    \textbf{Copies:} #7\bigskip\par
    \textbf{Numbered Pages:} \pageref{LastPage}\bigskip\par
    \textbf{Date of Completion:}\\ #8
  }
}

%Create danish project info
\newcommand{\danishprojectinfo}[8]{%
  \parbox[t]{\titlepageleftcolumnwidth}{
    \textbf{Titel:}\\ #1\bigskip\par
    \textbf{Tema:}\\ #2\bigskip\par
    \textbf{Projektperiode:}\\ #3\bigskip\par
    \textbf{Projektgruppe:}\\ #4\bigskip\par
    \textbf{Deltager(e):}\\ #5\bigskip\par
    \textbf{Vejleder(e):}\\ #6\bigskip\par
    \textbf{Oplagstal:} #7\bigskip\par
    \textbf{Sidetal:}  \pageref{LastPage} \bigskip\par
    \textbf{Afleveringsdato:}\\ #8
  }
}



% Preamblet indeholder alle de ting, som skal angives før selve indholdet af
% dokumentet starter. Kommandoen \input svarer til at kopiere indholdet fra
% filen ind, der hvor kommandoen kaldes.
% LaTeX skabelon, preamble-fil
% ------------------------------------------------------------------------------
% Dette er dokumentets preamble, som inkluderer ekstra pakker og definerer
% ekstra kommandoer.


% Basispakker ------------------------------------------------------------------

\usepackage[T1]{fontenc}     % understøttelse til at vise specialtegn i pdf'en
\usepackage[utf8]{inputenc}  % understøttelse til at inputte specialtegn
\usepackage[danish]{babel}   % oversættelse af automatiske ting
\usepackage[square]{natbib}  % sæt referencer til at bruger firkantklammer
\usepackage{hyperref}        % indsæt klik-bare hyperlinks
\usepackage{lastpage}        % referer til sidste side; brugt til at tælle sider


% Matematikpakker --------------------------------------------------------------

\usepackage{mathtools}  % matematik-kommandoer og -environments
\usepackage{amssymb}    % yderligere matematiksymboler
\usepackage{amsthm}     % environments til definitioner, sætninger, osv.
\usepackage{bbm}        % bedre 'blackboard bold' styling


% Figurpakker ------------------------------------------------------------------

\usepackage[dvipsnames]{xcolor}    % til farvet tekst
\usepackage{graphicx}              % håndtering af billedfiler
\usepackage{booktabs}              % bedre tabeller
\usepackage{tikz}                  % til at tegne vektorgrafik
\usepackage{algorithm}             % algoritme-floats
\usepackage[noend]{algpseudocode}  % pseudokode
\usepackage{listings}              % source code


% Layoutpakker mm. -------------------------------------------------------------

\usepackage{microtype}             % bedre mellemrum og font kerning
\usepackage{emptypage}             % fjern sidetal fra tomme sider
\usepackage{fancyhdr}              % headers and footers
\usepackage{titlesec}              % overskrifter
\usepackage[margin=3cm]{geometry}  % justering af sidemargen
\usepackage[toc,page]{appendix}    % bedre appendiks


% Skrifttype -------------------------------------------------------------------
% Skrifttyper er med pdflatex begrænset til dem, der findes pakker til (i
% modsætning til f.eks. xelatex, som lader dig bruge system fonts). Dokumentets
% skrifttype ændres ved blot at importere en pakke til det.
%
% Se
% * http://www.tug.org/pracjourn/2006-1/hartke/hartke.pdf
% * http://www.tug.dk/FontCatalogue/

\usepackage{lmodern}  % bedre standardfont


% Kommandoer mm. ---------------------------------------------------------------

% Danske oversættelser, som ikke fanges af babel-pakken
\renewcommand\appendixpagename{Appendicer}
\renewcommand\appendixtocname{Appendicer}
\renewcommand\lstlistingname{Script}
\renewcommand\lstlistlistingname{Scripts}
\makeatletter
\renewcommand\ALG@name{Algoritme}
\makeatother

% Theorem enviroments
\theoremstyle{plain}
\newtheorem{thm}{Sætning}[chapter]
\newtheorem{lem}[thm]{Lemma}
\newtheorem{prop}[thm]{Proposition}
\newtheorem*{cor}{Korollar}

\theoremstyle{definition}
\newtheorem{defn}[thm]{Definition}
\newtheorem{exmp}[thm]{Eksempel}

% Talmængder
\newcommand{\N}{\mathbb{N}}  % naturlige tal
\newcommand{\Z}{\mathbb{Z}}  % heltal
\newcommand{\Q}{\mathbb{Q}}  % rationelle tal
\newcommand{\R}{\mathbb{R}}  % reelle tal
\newcommand{\C}{\mathbb{C}}  % komplekse tal

% Andre operatorer og begreb
\newcommand{\1}{\mathbbm{1}}     % indikatorfunktion
\newcommand{\bigO}{\mathcal{O}}  % store O

% Headers og footers
\fancyhead{}                      % ryd standardindhold i header
\fancyhead[LE,RO]{\projectgroup}  % venstre på lige sider, højre på ulige
\fancyhead[RE]{\leftmark}         % kapitelnavn
\fancyhead[LO]{\rightmark}        % sektionsnavn
\fancyfoot{}                      % ryd standardindhold i footer
\fancyfoot[CE,CO]{\thepage}       % centrerede sidetal
\pagestyle{fancy}                 % aktiver 'fancy' style

% Titler/overskrifter
% http://mirrors.dotsrc.org/ctan/macros/latex/contrib/titlesec/titlesec.pdf
\titleformat
{\chapter}
[hang]
{\Huge\bfseries}
{\thechapter\enspace\textcolor{BrickRed}{|}\enspace}
{0pt}
{\Huge\bfseries}

% Foruddefinerede TikZ node styles
\tikzstyle{point} = [fill,shape=circle,minimum size=3pt,inner sep=0pt]
\tikzstyle{edge} = [fill=white,midway,inner sep=1pt]

% Syntax highlighting for Python-filer
\lstdefinestyle{custompy}{
  belowcaptionskip=\baselineskip,
  breaklines=true,
  language=Python,
  showstringspaces=false,
  basicstyle=\small\ttfamily,
  keywordstyle=\bfseries\color{green!40!black},
  commentstyle=\itshape\color{purple!40!black},
  identifierstyle=\color{blue},
  stringstyle=\color{orange},
}
\lstset{language=Python,style=custompy,captionpos=b}

% Alt tekst skal inputtes i et 'document' environment.
\begin{document}

% \frontmatter angiver, at det efterfølgende indhold ikke skal tælles med som
% nummererede sider, men istedet have deres egne sidetal (i romertal som
% standard).
\frontmatter
% LaTex skabelon, eksempel på forside
% ------------------------------------------------------------------------------

\begin{titlepage}
  \centering

  % AAU logo
  \includegraphics[width=\textwidth]{fig/img/aau-logo.pdf}

  \textsc{\LARGE P\projectnumber{}-Projekt}\\
  \textsc{\Large \projectdegree}

  \vspace{2cm}

  % Title
  {\Huge\bfseries\projecttitle}

  \vspace{0.5cm}

  {\Large\bfseries\sffamily\projectsubtitle}

  \vspace{2cm}

  % Authors and supervisors
  \begin{minipage}[t]{0.4\textwidth}
    \begin{flushleft}
      \large
      \emph{Forfattere:}\\
      \projectauthors
    \end{flushleft}
  \end{minipage}
  \begin{minipage}[t]{0.4\textwidth}
    \begin{flushright}
      \large
      \emph{Vejledere:}\\
      \projectsupervisors
    \end{flushright}
  \end{minipage}

  \vfill

  {\large\today}
\end{titlepage}

{\pdfbookmark[0]{English title page}{label:titlepage_us}
  \aautitlepage{%
    \englishprojectinfo{
      \LaTeX{} Template  % project title
    }{%
      Project Management  % project theme
    }{%
      Fall Semester 2018  % project period
    }{%
      Example Group  % project group
    }{
      % list of group members
      Nova B. Valberg-Madsen\\
      Line B. Olsen
    }{%
      % list of supervisors
      Janus S. Valberg-Madsen
    }{%
      0  % number of printed copies
    }{%
      \today  % date of completion
    }%
  }{%
    % Department and address
    \textbf{Dept. of Mathematical Sciences}\\
    Skjernvej 4A\\
    DK-9220 Aalborg Ø\\
    \href{http://math.aau.dk}{http://math.aau.dk}
  }{
    This document serves as a \LaTeX{} project template for first year students at Aalborg University, Dept. of Mathematics.
  } 







% Automatisk indholdsfortegnelse baseret på overskrifterne i inputfilerne.
\tableofcontents

% \mainmatter angiver hovedindholdet, som nummereres med arabiske tal.
\mainmatter

% Inputfiler bør opdeles, således at hver fil indeholdet ét kapitel.
% Kommandoen \include indsætter en ny side efterfulgt af indholdet i inputfilen.
% LaTeX template, example input file
% ------------------------------------------------------------------------------

\chapter{First Example}
\label{ch:first-example}

This is an example chapter with content.
Aside from \verb!\chapter!, there are several levels available to use for partitioning your body text:
%
\begin{itemize}
\item \texttt{part}
\item \texttt{chapter}
\item \texttt{section}
\item \texttt{subsection}
\item \texttt{subsubsection}
\item \texttt{paragraph}
\item \texttt{subparagraph}
\end{itemize}
%
Each level is a subsection of the above level.
Titles are added automatically to the table of contents.
See more at \url{https://en.wikibooks.org/wiki/LaTeX/Document_Structure#Sectioning_commands}.


% ------------------------------------------------------------------------------
\section{Environments}
\label{sec:sec-within-chap}

In \LaTeX{}, you are going to be using many different kinds of \emph{environments}.
These are scopes denoted with \verb!\begin{...}! and \verb!\end{...}!, enclosing special content such as lists, figures, equations, etc.
Table \ref{tab:my-table} lists some commonly used environments.

% LaTeX skabelon, eksempel på inputfil med en tabel
% ------------------------------------------------------------------------------
% Denne fil indeholder LaTeX-kode til at indsætte en enkelt tabel.  Bemærk at
% denne kode godt kan indsættes direkte i hovedfilerne, men jeg kan godt lige at
% opdele mit indhold i flere filer.


\begin{table}[htbp]
  \centering
  \begin{tabular}{ll}
    \toprule
    \textbf{Environment} & \textbf{Formål} \\
    \midrule
    \texttt{document} & Selve dokumentets indhold \\
    \texttt{table} & Tabeller ligesom denne \\
    \texttt{figure} & Figurer \\
    \texttt{equation} & Nummererede ligninger \\
    \texttt{align} & Justerede ligninger \\
    \texttt{itemize} & Punktliste \\
    \texttt{enumerate} & Nummereret liste \\
    \texttt{description} & Liste med beskrivelser \\
    \bottomrule
  \end{tabular}
  \caption{Almindelige \LaTeX{} environments og deres formål}
  \label{tab:my-table}
\end{table}


\subsection{Lists}
\label{sec:lists}

There are three essential list structures: \texttt{itemize}, \texttt{enumerate}, and \texttt{description}.
The \texttt{itemize} variant produces a simple bullet list.
Each item in the list are prepended by the \verb!\item! command.
%
\begin{verbatim}
  \begin{itemize}
    \item First item
    \item Second item
    \item Third item
  \end{itemize}
\end{verbatim}
%
\begin{itemize}
\item First item
\item Second item
\item Third item
\end{itemize}
%
The \texttt{enumerate} variant uses the same syntax for items as \texttt{itemize}, but produces a numbered list.
%
\begin{verbatim}
  \begin{enumerate}
    \item First item
    \item Second item
    \item Third item
  \end{enumerate}
\end{verbatim}
%
\begin{enumerate}
\item First item
\item Second item
\item Third item
\end{enumerate}
%
Finally, the \texttt{description} list in which \verb!\item! is given an item name as an optional argument, and the contents of the line is a description of that item.
This produces a list where the item names are typeset in bold followed by their descriptions as normal text.
%
\begin{verbatim}
  \begin{description}
    \item[First item] Description of first item
    \item[Second item] Description of second item
    \item[Third item] Description of third item
  \end{description}
\end{verbatim}
%
\begin{description}
\item[First item] Description of first item
\item[Second item] Description of second item
\item[Third item] Description of third item
\end{description}


\subsection{Equations}
\label{sec:equations}

One of the main reasons why people use \LaTeX{} is the beautiful math typesetting.
There are several different math environments to suit your needs, and most come in a numbered and unnumbered variants.
For example, the code
%
\begin{verbatim}
\begin{equation}
  \label{eq:1}
  e^{i\pi} - 1 = 0
\end{equation}
\end{verbatim}
%
produces the ouput
%
\begin{equation}
  \label{eq:1}
  e^{i\pi} + 1 = 0 ,
\end{equation}
%
and since it was given a label, it can be referenced with the command \verb!\eqref{eq:1}!, which produces a clickable reference in parentheses, \eqref{eq:1}.
If instead of \texttt{equation} you put \texttt{equation*}, the equation does not get a number.
Equivalently, you can use \verb!\[...\]!, so the code \verb!\[ e^{i\pi} + 1 = 0 \]! produces
\[ e^{i\pi} + 1 = 0 .\]

If you need multiple, aligned equations, e.g. for step-by-step calculations, use the \texttt{align} environment, which aligns the contents at \texttt{\&} characters.
For example,
%
\begin{verbatim}
\begin{align*}
  (x + y)^{2} &= x^{2} + xy + yx + y^{2} \\
              &= x^{2} + y^{2} + 2xy
\end{align*}
\end{verbatim}
%
produces
%
\begin{align*}
  (x + y)^{2} &= x^{2} + xy + yx + y^{2} \\
              &= x^{2} + y^{2} + 2xy .
\end{align*}
%
The double backslash denotes a line break.
Note the asterisk; like with \texttt{equation}, \texttt{align} has both a numbered and unnumbered version.
The numbered version has a seperate number for each line.

For rendering inline math, e.g. \(\cos^{2}\theta + \sin^{2}\theta = 1\), use \verb!\(...\)!.
Alternatively, you can also use \verb!$...$!, but \verb!\(...\)! has improved spacing and error messages.
%
See \url{https://en.wikibooks.org/wiki/LaTeX/Mathematics} for a good reference of symbols and commands.


\subsection{Figures}
\label{sec:figures}

Figures in \LaTeX{} are input as so-called \emph{floats} using the \texttt{figure} environment.
A floating object cannot be broken over a page, so the figure will be repositioned depending on the available space on the page.
The syntax is as follows:
%
\begin{verbatim}
\begin{figure}[placement]
  \centering
  \includegraphics[options]{path/to/image}
  \caption{The figure caption}
  \label{fig:label}
\end{figure}
\end{verbatim}
%
The optional argument \texttt{placement} can be either of \texttt{h} (here), \texttt{t} (top of page), \texttt{b} (bottom of page), or \texttt{p} (put on special page with only floats).
The \verb!\centering! command is there to center the image.
Among the options available for \verb!\includegraphics!, the most important one for you will probably be \texttt{width}.
To make the image take up half the page (within margins), use \verb!width=0.5\textwidth!.
All figures should have a caption, which is set with the \verb!\caption! command, and the \verb!\label! lets us reference it (for example, \verb!Figure \ref{fig:me}! becomes Figure \ref{fig:me}).
For more info, see \url{https://en.wikibooks.org/wiki/LaTeX/Floats,_Figures_and_Captions}.

% LaTeX template, example figure input file
% ------------------------------------------------------------------------------
% This file contains all the code involved in inputting a single image figure.
% While this code can just as well be inserted directly into the documents with
% the body text, I like to keep my projects as modular as possible.


\begin{figure}[htbp]
  \centering
  \includegraphics[width=0.5\textwidth]{fig/img/me.jpg}
  \caption{A picture of me responding to emails from my students}
  \label{fig:me}
\end{figure}


Figures need not necessarily consist of images like JPG, PNG, or PDF files, but can also consist of \LaTeX{} code.
An important example of this is \emph{TikZ}, a native vector graphics framework in \LaTeX{}.
A TikZ picture is enclosed in the environment \texttt{tikzpicture}, which can be enclosed in a \texttt{figure} environment, just like images:
%
\begin{verbatim}
\begin{figure}[placement]
  \centering
  \begin{tikzpicture}
    % tikz code goes here ...
  \end{tikzpicture}
  \caption{The figure caption}
  \label{fig:label}
\end{figure}
\end{verbatim}
%
An example of a graph drawn with TikZ code is shown on Figure \ref{fig:graph}.

% LaTeX skabelon, eksempel på inputfil med tikzpicture
% ------------------------------------------------------------------------------
% Denne fil indeholder LaTeX-kode til at indsætte en enkelt figur med TiKZ.
% Bemærk at denne kode godt kan indsættes direkte i hovedfilerne, men jeg kan
% godt lige at opdele mit indhold i flere filer.


\begin{figure}[htbp]
  \centering

  \begin{tikzpicture}

    % Nodes
    \node[point] at (0,1) (a) [label=left:\(a\)] {};
    \node[point] at (1,2) (b) [label=above:\(b\)] {};
    \node[point] at (1,0) (c) [label=below:\(c\)] {};
    \node[point] at (3,2) (d) [label=above:\(d\)] {};
    \node[point] at (3,0) (e) [label=below:\(e\)] {};
    \node[point] at (4,1) (z) [label=right:\(z\)] {};

    % Edges
    \footnotesize
    \draw (a) -- (b) node[edge] {4};
    \draw (a) -- (c) node[edge] {2};
    \draw (b) -- (c) node[edge] {1};
    \draw (b) -- (d) node[edge] {5};
    \draw (c) -- (d) node[edge] {8};
    \draw (c) -- (e) node[edge] {10};
    \draw (d) -- (e) node[edge] {2};
    \draw (d) -- (z) node[edge] {6};
    \draw (e) -- (z) node[edge] {3};
    
  \end{tikzpicture}

  \caption{Graf fra \citep{rosen}.}
  \label{fig:graph}
\end{figure}


To learn TikZ syntax, see \citep{minimaltikz}, the tutorial sections of \citep{tikzmanual}, and \url{https://en.wikibooks.org/wiki/LaTeX/PGF/TikZ}.
% LaTeX template, example input file
% --------------------------------------------------------------------------------------------------

\chapter{Second Example}
\label{ch:second-example}

Here is another example input file.

\section{Custom Environments and Commands}
\label{sec:custom}

While \LaTeX{} provides commands for many different purposes, you will often find yourself defining your own.
For this template, I have included some examples of custom environments and commands in the preamble (\texttt{premable.tex}).

Such commands can save you a lot of typing when working on a long, modular document such as a semester project.
For instance, instead of typing \verb!\mathbb{N}! every time you want the symbol for the set of natural numbers, define a shorter command, like \verb!\N!.
The syntax for defining commands is as follows:
%
\begin{verbatim}
\newcommand{name}[num]{definition}
\end{verbatim}
%
where \texttt{name} is the command name, e.g. \verb!\N!, \texttt{num} is the number of arguments the command takes (omit the square brackets if the command takes no arguments), and \texttt{definition} is the output of the command, e.g. \verb!\mathbb{N}!.


\subsection{Definitions, Theorems, Proofs}
\label{sec:thms}

In a mathematics project, you are going to be including mathematical definitions, propositions, lemmas, theorems, etc.
The \texttt{amsthm} package provides a simple way to define such environments:
%
\begin{verbatim}
\newtheorem{name}{Printed output}[numberby]
\end{verbatim}
%
A few examples are included in the preamble.
See \url{https://en.wikibooks.org/wiki/LaTeX/Theorems} for more information.
%
\begin{defn}[Rational numbers]
  A real number \(r\) is called \emph{rational} if there exist integers \(p\) and \(q\) with \(q \neq 0\) such that \(r = p/q\)
  (a real number which is not rational is called \emph{irrational}).
  The set of rational numbers is denoted with \(\Q\).
\end{defn}
%
\begin{thm}[Example of Named Theorem]
  \(\sqrt{2}\) is irrational.
\end{thm}
%
\begin{proof}
  Suppose \(\sqrt{2} \in \Q\) and let \(k = \min\{\, N \in \N : N\sqrt{2} \in \N \,\}\).
  But then \(k (\sqrt{2} - 1) = (k\sqrt{2} - k) \in \N\), and \(k(\sqrt{2} - 1) < k\), which is a contradiction.
  Therefore, \(\sqrt{2} \notin \Q\).
\end{proof}


\section{Citations}
\label{sec:citations}

When you use other people's work, you must cite them.
Citations in \LaTeX{} is usually handled with BibTeX, a procedure that identifies resources from a bibliography file you provide.
Literature resources must be defined in a \texttt{.bib} file and included with the \verb!bibliography! command.
Each entry in this file must have a resource type and a key to identify it by.
The syntax for citing a resource is
%
\begin{verbatim}
\citep[scope]{bibkey}
\end{verbatim}
%
where \texttt{bibkey} is the unique identifier you give the resource in the \texttt{.bib} file, and \texttt{scope} can be used to refer to a specific chapter or range of pages.

When you cite a resource, that resource is automatically added to the literature list, and in the PDF file the citations become clickable links pointing to that list.
For example, by citing \citep[pp. 104-110]{edpenn}, an entry for it is added to the bibliography on the last page.

There are several different types of BibTeX entries, and each type has its own set of mandatory and optional fields.
For more information, see \url{https://en.wikibooks.org/wiki/LaTeX/Bibliography_Management#BibTeX}.

% Appendicer
\begin{appendices}
  % LaTeX skabelon, eksempel på appendiks-fil
% ------------------------------------------------------------------------------

\chapter{Pseudokode}
\label{app:pseudocode}

For at indsætte pseudokode i ens dokument bruges \texttt{algorithmic} environmentet, som placeres inde i et \texttt{algorithm} environment for at få en figur-agtig float:
%
\begin{verbatim}
\begin{algorithm}
  \begin{algorithmic}
    % pseudokode her...
  \end{algorithmic}
  \caption{algoritmens titel}
  \label{alg:label}
\end{algorithm}
\end{verbatim}

% LaTeX skabelon, eksempel på inputfil med pseudokode
% ------------------------------------------------------------------------------
% Denne fil indeholder al LaTeX-kode, der skal bruges til at indsætte en enkel
% pseudokode-blok.  Bemærk at denne kode godt kan indsættes direkte i
% hovedfilerne, men jeg kan godt lige at opdele mit indhold i flere filer.


\begin{algorithm}[htbp]
  \begin{algorithmic}
    \Procedure{bubblesort}{\,\(a_{1}, \dotsc, a_{n}\) : reelle tal med \(n \geq 2\)\,}
    \For{\(i \coloneqq 1\) to \(n - 1\)}
    \For{\(j \coloneqq 1\) to \(n - i\)}
    \If{\(a_{j} > a_{j+1}\)}
    swap \(a_{j}\) and \(a_{j+1}\)
    \EndIf
    \EndFor
    \EndFor
    \EndProcedure
    \Comment{\(a_{1}, \dotsc, a_{n}\) er nu i voksende orden}
  \end{algorithmic}
  \caption{Bubble Sort}
  \label{alg:some}
\end{algorithm}


Se \url{https://en.wikibooks.org/wiki/LaTeX/Algorithms#Typesetting_using_the_algorithmicx_package} for en oversigt over syntaksen.
  % LaTeX skabelon, eksempel på appendiks-fil
% ------------------------------------------------------------------------------

\chapter{Python Kode}
\label{app:code}

Hvis I implementerer algoritmer i et specifikt programmeringssprog, vil I muligvis inkluderer kildekoden i et appendiks som dette.
Pakken \texttt{listings} gør det muligt at inkludere kildekode med syntax highlighting for mange forskellige sprog.

Med environmentet \texttt{lstlisting} kan man indsætte kode direkte, men for det meste vil I typisk bruge \verb!\lstinputlisting! til at indsætte kode fra en fil, f.eks.
%
\begin{verbatim}
\lstinputlisting[options]{path/to/source_file}
\end{verbatim}
%
Nogle nyttige \texttt{options} for denne kommando er:
%
\begin{description}
\item[caption] sætter caption for koden
\item[label] sætter label for koden (til \verb!\ref!)
\item[language] sætter syntaksen til et bestemt sprog lokalt
\item[style] sætter highlighter stilen lokalt
\item[firstline] springer linjerne før den angivede linje over 
\item[lastline] stopper med at læse linjer, når den angivede linje nåes
\end{description}
%
I preamblet kan man også sætte globale options med \verb!\lstset!.
For eksempel, hvis al ens kildekode er skrevet i Python, kan man skrive \verb!\lstset{language=Python}!, og så behøver man ikke bruge sætte \texttt{language} i hver enkel \verb!\lstinputlisting!.

\lstinputlisting[firstline=5,caption=Bubble Sort in Python,label=py:bubble]{py/bubblesort.py}

For mere information, se \url{https://en.wikibooks.org/wiki/LaTeX/Source_Code_Listings}.
\end{appendices}

% \backmatter er til ekstra ting og sager. Overskrifter er ikke nummererede.
\backmatter

% Automatiske lister af alle figurer og tabeller, baseret på deres captions.
\listoffigures
\listoftables

% Automatisk liste af referencer, baseret på, hvilke referencer fra bib-filen,
% der rent faktisk blev henvist til i dokumentet.
\bibliographystyle{apalike}
\bibliography{incl/bib/bibliography}

\end{document}
